% Generated draft: requires mathematical and source review.
\documentclass[11pt]{article}
\usepackage[margin=25mm]{geometry}
\usepackage{fontspec}
\setmainfont{Latin Modern Roman}
\usepackage{amsmath,amssymb,mathtools,unicode-math}
\setmathfont{Latin Modern Math}
\usepackage{enumitem,needspace,xurl,hyperref}
\setlength{\parindent}{0pt}
\setlength{\parskip}{5pt}
\setlength{\emergencystretch}{3em}
\newcommand{\N}{\mathbb{N}}
\newcommand{\Z}{\mathbb{Z}}
\newcommand{\Q}{\mathbb{Q}}
\newcommand{\R}{\mathbb{R}}
\newcommand{\C}{\mathbb{C}}
\newcommand{\F}{\mathbb{F}}
\newcommand{\A}{\mathbb{A}}
\newcommand{\PP}{\mathbb{P}}
\newcommand{\E}{\mathbb{E}}
\newcommand{\eps}{\varepsilon}
\newcommand{\dd}{\,\mathrm{d}}
\begin{document}
\section*{UnsolvedMath ID: 800871}
\textbf{Source problem code:} catalogue entry
\subsection{Definitions and mathematical statement}
Let \(n\geq 1\) be an integer and let \(\Omega\subset\C^n\) be a domain. A function is plurisubharmonic if it is upper semicontinuous and its restriction to each complex line is subharmonic wherever defined. Let \(\varphi\) and \(\varphi_j\), \(j\geq 1\), be bounded plurisubharmonic functions on \(\Omega\), with
\[
\varphi_j\longrightarrow\varphi\quad\text{in }L^1_{\mathrm{loc}}(\Omega).
\]
Thus \(\int_K|\varphi_j-\varphi|\,\dd V\to0\) for every compact \(K\subset\Omega\), where \(\dd V\) is Lebesgue measure. Each function is bounded individually; the supplied statement does not impose a bound uniform in \(j\).

For definiteness, use \(d^c=\frac{i}{2}(\bar\partial-\partial)\), so \(dd^c=i\partial\bar\partial\). For a locally bounded plurisubharmonic function \(u\), write \((dd^c u)^n\) for its positive Bedford--Taylor complex Monge--Amp\`ere measure, defined by the locally bounded pluripotential product. The normalization of \(d^c\) is not specified in the source; the convention here fixes the scale of these measures.

To give a precise interpretation of cluster points, set \(\mu_j=(dd^c\varphi_j)^n\) and define
\[
\mathcal C((\varphi_j))=
\left\{\mu\in\mathcal M_+(\Omega):
\begin{array}{l}
\text{there exist }j_1<j_2<\cdots\text{ such that}\\
\displaystyle\lim_{k\to\infty}\int_\Omega f\,\dd\mu_{j_k}
=\int_\Omega f\,\dd\mu
\quad\text{for every }f\in C_c(\Omega)
\end{array}\right\},
\]
where \(\mathcal M_+(\Omega)\) denotes the positive Radon measures on \(\Omega\), and \(C_c(\Omega)\) denotes the continuous compactly supported test functions. This uses local weak convergence of measures; the supplied source does not explicitly specify a topology.

\textbf{Question (attributed to Bedford).} Given \(\varphi\) as above, what can be said about \(\mathcal C((\varphi_j))\) for sequences satisfying these hypotheses? In particular, describe the possible cluster measures under the stated convergence assumption.

The source asks an open-ended question and supplies no proposed characterization or particular property to establish. It does not explicitly distinguish the cluster set for one chosen sequence from the collection of limits obtainable as the approximating sequence varies. No boundary conditions or existence of a cluster point are asserted.
\subsection{Short English statement}
For bounded plurisubharmonic approximations converging locally in \(L^1\) to a bounded plurisubharmonic function, describe the possible subsequential limits of their complex Monge--Amp\`ere measures. The draft interprets limits through compactly supported continuous test functions. The source leaves the convergence topology and the intended form of a characterization unspecified.
\subsection{Sources}
\begin{itemize}
\item Local UnsolvedMath registry, permanent ID \(800871\), source code \texttt{catalogue entry}, title: \emph{Open Problems in Pluripotential Theory --- Question 2}. The supplied record labels the problem open; this is a source label, not an independently verified current assessment.
\item Dinew, Guedj, and Zeriahi, \emph{Open Problems in Pluripotential Theory} (2015), Question 2. Supplied source URL: \url{https://arxiv.org/abs/1511.00705}. The question is explicitly attributed to Bedford in the supplied statement; the registry lists Dinew, Guedj, and Zeriahi as proposers.
\item Supplied provenance: extraction method \texttt{arxiv\_latex}; source access date \texttt{2026-10-16}. This access date and provenance are supplied metadata, not a claim of independent source consultation or verification.
\end{itemize}
\end{document}
